We present FTRFS, a Linux filesystem for radiation-robust
embedded environments. Building on the original concept of Fuchs,
Langer and Trinitis (ARCS 2015), the present implementation is a
contemporary response to the 2015 problem statement under the
standards and tooling that have emerged in the intervening eleven
years: the Linux kernel's mature \texttt{lib/reed\_solomon}
infrastructure, MIL-STD-882E hazard tracking practice for
software-intensive systems, and reproducible embedded builds via
the Yocto Project.

FTRFS combines universal Reed--Solomon forward error correction on
metadata (inodes, allocation bitmap, superblock), a persistent
on-disk radiation event journal, and fail-closed mount semantics.
To our knowledge, no other Linux filesystem combines these three
properties.

This Technical Report v1 documents the project state at commit
\texttt{9a63468}: stage 3 partial closure, format v3 stable, format
v4 design frozen but not yet implemented, validation on a hardened
arm64 build (Yocto Scarthgap, linux-mainline 7.0.1, zero unpatched
kernel CVE), and an I/O baseline captured on a four-node libvirt
cluster. The project's main upstream review thread is publicly
archived on lore.kernel.org~\cite{ftrfs-rfc-thread}.

This is part of a versioned publication roadmap (v2--v4) following
the project's engineering milestones.
